% Chapter Template
% Save as chapters/ch01-name.tex

\chapter{Chapter Title}
\label{ch:chapterlabel}

Opening paragraph introducing the chapter content. This should set context
and explain what the reader will learn.

% ============================================================================
\section{First Major Section}
% ============================================================================

Section introduction text.

\subsection{Subsection Title}

Regular paragraph text with \textbf{bold} and \textit{italic} formatting.
Code inline looks like \texttt{variable\_name} or \texttt{function()}.

% Problem-Pattern structure (common in technical guides)
\begin{problembox}
Describe the problem or challenge that readers face. Be specific about
the symptoms and why this matters.
\end{problembox}

\begin{solutionbox}
Present the solution or pattern that addresses the problem. Include
specific, actionable guidance.
\end{solutionbox}

\subsection{Code Examples}

Example with syntax highlighting:

\begin{lstlisting}[language=typescript]
interface User {
  id: string;
  name: string;
  email: string;
}

async function getUser(id: string): Promise<User> {
  const response = await fetch(`/api/users/${id}`);
  return response.json();
}
\end{lstlisting}

Bash commands:

\begin{lstlisting}[language=bash]
# Install dependencies
npm install

# Run tests
npm test

# Build for production
npm run build
\end{lstlisting}

\subsection{Tables}

\begin{table}[htbp]
\centering
\small
\begin{tabularx}{\textwidth}{lXl}
\toprule
\textbf{Column 1} & \textbf{Column 2 (Flexible)} & \textbf{Column 3} \\
\midrule
Item A & Description that can wrap to multiple lines if needed & Value \\
Item B & Another description & Value \\
Item C & Yet another description & Value \\
\bottomrule
\end{tabularx}
\caption{Descriptive caption for the table}
\label{tab:tablename}
\end{table}

Reference the table: See Table~\ref{tab:tablename}.

\subsection{Figures and Diagrams}

\begin{figure}[htbp]
\centering
\begin{tikzpicture}[
    node distance=2cm,
    box/.style={draw, rounded corners, minimum width=2cm, minimum height=1cm, align=center}
]
    \node[box, fill=blue!10] (a) {Step 1};
    \node[box, fill=orange!10, right=of a] (b) {Step 2};
    \node[box, fill=green!10, right=of b] (c) {Step 3};

    \draw[->, thick] (a) -- (b);
    \draw[->, thick] (b) -- (c);
\end{tikzpicture}
\caption{Workflow diagram showing the three-step process}
\label{fig:workflow}
\end{figure}

Reference the figure: See Figure~\ref{fig:workflow}.

\subsection{Lists}

Bullet list:
\begin{itemize}
    \item First item
    \item Second item with more detail
    \item Third item
\end{itemize}

Numbered list:
\begin{enumerate}
    \item First step
    \item Second step
    \item Third step
\end{enumerate}

% ============================================================================
\section{Second Major Section}
% ============================================================================

Continue with more content...

\begin{tipbox}
Helpful tips or best practices go in tip boxes. These highlight
information the reader should pay special attention to.
\end{tipbox}

\begin{warningbox}
Warnings about common mistakes or potential issues. Use sparingly
for genuine cautions.
\end{warningbox}

% ============================================================================
\section{Summary}
% ============================================================================

Concluding section that recaps key points:

\begin{table}[htbp]
\centering
\begin{tabularx}{\textwidth}{lX}
\toprule
\textbf{Concept} & \textbf{Key Takeaway} \\
\midrule
Concept 1 & Main point to remember \\
Concept 2 & Another key insight \\
Concept 3 & Final important point \\
\bottomrule
\end{tabularx}
\end{table}

% ============================================================================
\section{Related}
% ============================================================================

Cross-references to other chapters:
\begin{itemize}
    \item Chapter~\ref{ch:another} --- Related topic
    \item Appendix~\ref{app:reference} --- Detailed reference material
\end{itemize}
